\section{The native zeta product and its natural boundary}\label{sec:zeta}

Define the dynamical zeta function for the same map and clock by
\begin{equation}\label{eq:zeta-def}
 Z_S(t)=\exp\left(\sum_{n\ge1}\frac{N_n}{n}t^n\right).
\end{equation}
No Euler factors from another object enter this definition.
For the applicable support class put
\begin{equation}\label{eq:ABtheta}
 (A,B_0,\theta)=
 \begin{cases}
 \displaystyle\left(\frac{\ell-1}{q-1},
       \frac{q(h+1)-2p}{q-1},\frac1q\right),&\text{high support},\\[8pt]
 \displaystyle\left(\frac{p}{q-p},
       \frac{q-2p}{q-p},\frac pq\right),&\text{low support}.
 \end{cases}
\end{equation}
We use $B_0$ here to distinguish this real constant from the
coefficient $B=\sigma(b)$ in the perfected-ring proof.
Theorem~\ref{thm:main} is equivalent to
\begin{equation}\label{eq:normalized-count}
 \frac{N_n}{q^{2n}}=A+B_0\theta^{p^{v_p(n)}}.
\end{equation}
In either class,
\begin{equation}\label{eq:positive-parameters}
 A,B_0>0,\qquad 0<\theta<1,\qquad A+B_0\theta=\frac{2p}{q}.
\end{equation}
In particular the normalization used below is $A+B_0\theta$,
not an assumed identity $A+B_0=1$.

\begin{theorem}[Product and meromorphic natural boundary]\label{thm:zeta}
For $|u|<1$,
\begin{equation}\label{eq:zeta-product}
 Z_S(u/q^2)=(1-u)^{-2p/q}
       \prod_{k\ge1}(1-u^{p^k})^{e_k},\qquad
 e_k=\frac{B_0}{p^k}
       \left(\theta^{p^{k-1}}-\theta^{p^k}\right)>0.
\end{equation}
The product is defined by logarithms equal to zero at $u=0$ and
converges locally uniformly in the unit disk. For every positive
integer $m$, the circle $|t|=q^{-2}$ is a meromorphic natural
boundary for $Z_S(t)^m$: no arc admits meromorphic continuation.
In particular $Z_S$ has radius of convergence exactly $q^{-2}$
and is transcendental over $\C(t)$.
\end{theorem}

\begin{proof}
For $w=p^{v_p(n)}$ the finite telescoping identity is
\begin{equation}\label{eq:telescoping}
 \theta^w=\theta+
 \sum_{k\ge1}\left(\theta^{p^k}-\theta^{p^{k-1}}\right)
               \mathbf1_{p^k\mid n}.
\end{equation}
For each fixed $n$ only finitely many summands occur. Interchanging
this sum with $\sum u^n/n$ is valid uniformly on $|u|\le r<1$:
the sum of the absolute coefficient differences over $k$ is $\theta$,
so an upper bound for the resulting double absolute sum is
$\theta\sum_{n\ge1}r^n/n<\infty$. Also
\[
 \sum_{\substack{n\ge1\\p^k\mid n}}\frac{u^n}{n}
 =-\frac1{p^k}\log(1-u^{p^k}).
\]
Substitution of \eqref{eq:telescoping} into
\eqref{eq:normalized-count} and \eqref{eq:zeta-def} proves
\eqref{eq:zeta-product}, with locally uniform logarithmic
convergence. The resulting holomorphic function has no zeros
in the unit disk, as is also immediate from \eqref{eq:zeta-def}.

Let $\xi$ be a primitive $p^r$th root of unity, where $r\ge1$.
The restriction $r\ge1$ ensures $\xi\ne1$, so the initial
factor $1-u$ does not vanish along $u=\rho\xi$ as
$\rho\uparrow1$. The finitely many factors with $k<r$ likewise
have nonzero limits. For $k\ge r$, $\xi^{p^k}=1$, and
\[
 \log|1-(\rho\xi)^{p^k}|=\log(1-\rho^{p^k}).
\]
For $0<\rho<1$,
\[
 0\le\frac{\log(1-\rho^{p^k})}{\log(1-\rho)}\le1,
 \qquad
 \lim_{\rho\uparrow1}
 \frac{\log(1-\rho^{p^k})}{\log(1-\rho)}=1.
\]
The bound $\sum e_k<\infty$ permits dominated convergence.
Writing $F(u)=Z_S(u/q^2)$ gives the radial logarithmic order
\begin{equation}\label{eq:radial-order}
 \lim_{\rho\uparrow1}
 \frac{\log|F(\rho\xi)|}{\log(1-\rho)}
 =\eta_r:=\sum_{k\ge r}e_k>0.
\end{equation}
These positive tails tend to zero; more explicitly,
\begin{equation}\label{eq:eta-bound}
 0<\eta_r\le
 \frac{B_0}{p^r}\sum_{k\ge r}
       (\theta^{p^{k-1}}-\theta^{p^k})
 =\frac{B_0}{p^r}\theta^{p^{r-1}}.
\end{equation}

Fix a positive integer $m$. For all sufficiently large $r$,
$0<m\eta_r<1$. A nonzero meromorphic germ at $\xi$ has the
form $(u-\xi)^j G(u)$ with $j\in\Z$ and $G$ holomorphic and
nonzero at $\xi$. Its radial logarithmic order is the integer
$j$. Equation~\eqref{eq:radial-order} makes that impossible
for $F^m$ at every primitive $p^r$th root with $r$ sufficiently
large. Such roots are dense on the unit circle, so no open arc
can admit meromorphic continuation of $F^m$.

The defining series is holomorphic in $|u|<1$ by
\eqref{eq:normalized-count}; the boundary just established
precludes a larger disk. An algebraic function has only finitely
many singular or branch points in the finite plane and therefore
has continuation across some arc of this circle. Thus $F$, and
equivalently $Z_S$, is not algebraic over $\C(t)$.
\end{proof}

This proof is the positive-tail mechanism already used in
C404~\cite{C404}. The contribution of the present degree calculation
is to determine which of the two exact triples in
\eqref{eq:ABtheta} governs every coefficient choice.
